\chapter{Herniated Disc}
\index{Herniated Disc}
\label{sec:Herniated Disc}

\section{Overview}
A herniated disc occurs when material from the soft inner nucleus of an intervertebral disc protrudes through a weakened or torn outer annulus. The protrusion may irritate or compress a nearby nerve root, producing radicular pain, numbness, or weakness. The lumbar spine is most often affected, followed by the cervical spine.
\section{Causes}
Disc herniation usually reflects age-related dehydration and degeneration combined with mechanical stress. Lifting while twisting, repetitive loading, an acute injury, or rarely a genetic connective-tissue disorder may contribute. Many disc protrusions are visible on imaging but cause no symptoms.
\section{Risk Factors}
\begin{itemize}[noitemsep]
  \item Age-related disc degeneration
  \item Repetitive lifting, twisting, or vibration
  \item Sedentary lifestyle and weak trunk musculature
  \item Obesity and cigarette smoking
  \item Previous back injury or family history of disc disease
\end{itemize}
\section{Signs and Symptoms}
Lumbar herniation may cause low-back pain with pain, tingling, numbness, or weakness extending into the buttock and leg (sciatica). Cervical herniation may cause neck pain with symptoms extending into the shoulder, arm, or hand. Symptoms may worsen with coughing, sneezing, or prolonged sitting. New bowel or bladder dysfunction, saddle numbness, or progressive severe weakness may indicate cauda equina or significant spinal-cord compression and requires emergency evaluation.
\section{Progression and Stages}
Symptoms may begin suddenly or gradually. Many improve over several weeks as inflammation settles and the protruded material shrinks or is resorbed. Persistent pain, worsening neurologic deficits, or recurrent episodes require reassessment and may lead to specialist or surgical evaluation.
\section{Complications}
\begin{itemize}[noitemsep]
  \item Persistent radicular pain or sensory loss
  \item Muscle weakness and impaired mobility
  \item Chronic pain and reduced function
  \item Cauda equina syndrome or spinal-cord compression, rarely
\end{itemize}
\section{Diagnosis}
Diagnosis is based on the history and neurologic examination, including strength, sensation, reflexes, gait, and provocative maneuvers. MRI is the preferred imaging test when symptoms persist, neurologic deficits are significant, or urgent structural assessment is needed. Imaging findings must be correlated with symptoms because asymptomatic disc protrusions are common. Electrodiagnostic testing may help when the diagnosis remains uncertain.
\section{Prevention}
\begin{itemize}[noitemsep]
  \item Use safe lifting technique and avoid twisting while lifting
  \item Maintain regular exercise and trunk strength
  \item Keep a healthy weight and avoid smoking
  \item Alternate prolonged sitting with movement and use good posture
\end{itemize}
\section{Treatment Options}
Initial management commonly includes continued activity as tolerated, education, heat or cold, physical therapy, and analgesic or anti-inflammatory medicines when appropriate. Selected patients may receive an epidural injection. Surgery, usually discectomy or microdiscectomy, is considered for persistent disabling radicular symptoms, progressive neurologic deficit, or cauda equina syndrome. Urgent surgical assessment is required for suspected cauda equina syndrome.
\section{Living with It}
Patients should follow the rehabilitation plan, gradually resume normal activities, avoid prolonged bed rest, and seek review if pain worsens or weakness develops. Emergency care is needed for new bladder or bowel dysfunction, saddle numbness, or rapidly progressive weakness.
\section{Outlook}
Most people improve with conservative treatment. Recovery depends on the location and severity of nerve compression, neurologic findings, general health, and response to rehabilitation. Surgery can provide faster relief for selected patients with persistent nerve compression but is not necessary for most cases.
